% 03_related_work.tex
\section{Related Work}
\label{sec:related}

\paragraph{Passive detection.} DetectGPT~\citep{mitchell2023detectgpt}
exploits log-probability curvature; GPTZero~\citep{tian2023gptzero}
combines perplexity and burstiness. Both read statistical traces in
existing text, with no controllable false-positive frontier and sharp
degradation under
paraphrase~\citep{krishna2023paraphrasing,sadasivan2023can}. They
cannot serve as the technical substrate for Article~50 because neither
regulator nor deployer can set the detection regime in advance.

\paragraph{Active watermarking.} Green--red-list
watermarking~\citep{kirchenbauer2023watermark} biases next-token
logits toward a hashed green list and detects via a $z$-test on
green-token frequency. The scheme is well-understood but requires
temperature sampling, degrades under paraphrase, and keeps the
detection-quality tradeoff inside a single bias magnitude with no
analytic inversion back to a regulatory parameter.
\citet{christ2024undetectable} establish the formal quality--security
frontier for undetectable watermarks; the result bounds what any
scheme can achieve but does not yield a deployable construction.
\citet{aaronson2023semstamp} and the SemStamp line embed signals in
semantic-level features and require LLM queries at detection time,
foreclosing lightweight third-party audit.

\paragraph{Entropy-adaptive schemes.}
% FLAG: lee2023sweet citation details (authors/venue/year) are guessed.
% FLAG: lu2024ewd citation is unverified.
SWEET~\citep{lee2023sweet} and EWD~\citep{lu2024ewd} restrict
watermarking to \emph{high-entropy} positions, arguing the scheme has
the most room to inject signal where the model is least committed.
Our scheme inverts this policy: MCL watermarks at \emph{low}-entropy
positions, trading on the intuition that near-deterministic tokens
survive meaning-preserving rewrites at higher rates. We formalize and
test this inversion (\S\ref{sec:experiments}).

\paragraph{Technical AI governance.} Article~50 of the EU AI
Act~\citep{euaiact2024} and the EU's draft Code of Practice on
Transparency~\citep{eucodepractice2025} require machine-readable
marking of AI-generated content; the OECD Hiroshima
Process~\citep{oecdhiroshima2024} provides a parallel international
framework. These instruments leave the detection regime
(FPR, quality floor, robustness) implicit. MCL's contribution, beyond
the watermarking scheme itself, is that it surfaces these parameters
analytically, producing a policy-to-configuration map auditable by
third parties (\S\ref{sec:discussion}).
